\section{Introduction}\label{sec:intro}

Neural theorem provers have achieved remarkable success in recent years, with
systems such as GPT-f~\cite{polu2020generative},
Baldur~\cite{first2023baldur}, HyperTree Proof Search~\cite{lample2022hypertree},
and DeepSeek-Prover-V2~\cite{ren2025deepseek} proving substantial fractions of
benchmark problems in formal mathematics. A recurring empirical observation is
that systems employing chain-of-thought (CoT) decomposition---generating proofs
step by step---outperform those that attempt to produce entire proofs in a single
pass. Despite the practical importance of this phenomenon, its mathematical
underpinnings remain poorly understood.

One natural avenue for explanation lies in the geometry of the loss landscape.
The loss landscape of a neural network, viewed as a function from parameter
space to the reals, encodes the difficulty of the optimization problem: flatter
minima are associated with better generalization via PAC-Bayes
bounds~\cite{foret2021sharpness}, mode connectivity reveals the basin structure
of the solution space~\cite{frankle2020linear, kuditipudi2019explaining}, and
Hessian spectral properties characterize local curvature and
trainability~\cite{ghorbani2019investigation}. Powerful tools for analyzing
these properties---filter-normalized
visualization~\cite{li2018visualizing}, permutation alignment~\cite{ainsworth2023git},
and stochastic Hessian estimation---are now well established.

However, these tools have been applied almost exclusively to classification and
generation tasks on standard benchmarks. No prior work, to our knowledge,
compares loss landscape geometry across different proof generation strategies in
neural theorem proving. This gap is notable because theorem proving possesses
structural features absent from typical machine learning tasks: the existence of
a verification oracle (proofs can be checked mechanically), the presence of
multiple valid proofs per theorem, and the hierarchical structure of proof
derivations.

In this paper, we provide the first systematic geometric comparison of loss
landscapes arising from CoT versus direct proof generation. We make the
following contributions:

\begin{itemize}[leftmargin=*,itemsep=0pt,topsep=4pt]
\item We formalize the decomposed and monolithic loss functions for sequence
  prediction and prove that, under a conditional independence assumption on
  intermediate steps, the Hessian spectral norm of the decomposed loss is bounded
  by the maximum per-step spectral norm (Theorem~\ref{thm:hessian-bound}).

\item We establish a superadditivity property of loss barriers under linear
  interpolation for decomposed objectives, providing a partial explanation for the
  higher mode connectivity barriers observed in CoT models
  (Proposition~\ref{prop:barrier-super}).

\item We present a controlled experimental study comparing five geometric
  properties of the loss landscape---Hessian top eigenvalue, Hessian trace,
  SAM-style sharpness, mode connectivity barriers, and perturbation
  sensitivity---for matched transformer models trained on propositional logic
  theorem proving.

\item We identify a novel ``flat but isolated'' phenomenon: CoT minima are
  individually flatter yet more separated from each other in parameter space
  than direct minima, challenging the common assumption that flatter landscapes
  imply better-connected solution spaces.
\end{itemize}

The paper is organized as follows. Section~\ref{sec:prelim} establishes
notation, definitions, and prerequisite results from loss landscape theory and
neural sequence prediction. Section~\ref{sec:main} contains our main
theoretical results and the complete experimental analysis.
Section~\ref{sec:discussion} discusses implications, limitations, and
connections to related work. Section~\ref{sec:conclusion} summarizes our
findings and outlines directions for future research.
